La variación trimestral del Producto Interno Bruto (PIB) muestra un comportamiento claramente diferenciado antes y después del año 2020. Durante el período comprendido entre 2015 y 2019 se observa un patrón estacional relativamente estable, caracterizado por un crecimiento importante en el cuarto trimestre de cada año y una desaceleración o variación negativa en el segundo trimestre en la mayoría de los períodos analizados. Este comportamiento refleja la existencia de ciclos económicos recurrentes asociados a la dinámica productiva nacional.
El año 2020 representa un punto de quiebre en la serie histórica debido a la crisis provocada por la pandemia de COVID-19. Como consecuencia de las restricciones sanitarias y la paralización parcial de las actividades económicas, el PIB registró una contracción significativa durante los dos primeros trimestres, alcanzando una variación mínima cercana al -9.73%. Posteriormente, la reapertura gradual de la economía permitió una rápida recuperación de la actividad productiva, reflejada en una variación positiva cercana al 9.98% durante el tercer trimestre, constituyendo el mayor crecimiento observado en todo el período de estudio.
A partir de 2022 la economía guatemalteca retomó un comportamiento más estable, registrando tasas de crecimiento moderadas y sin variaciones trimestrales negativas significativas. Asimismo, vuelve a apreciarse el patrón estacional observado antes de la pandemia, lo que sugiere un proceso de normalización de la actividad económica tras el choque extraordinario experimentado en 2020. En conjunto, estos resultados evidencian la resiliencia de la economía guatemalteca y su capacidad para recuperar la senda de crecimiento una vez superado el período de mayor incertidumbre.